\section{What the results say}
\paragraph{Discrimination is not the bottleneck on short imbalanced text.} Eight model families spanning five decades of machine learning, from naive Bayes to a bidirectional LSTM, rank hate-speech tweets within 0.013 ROC-AUC of one another. What separates them in F1 is the calibration of their scores and whether a threshold chosen on fold models survives refitting. A study that compares such models at a fixed threshold is mostly comparing calibration, and a study that reports accuracy on a 93/7 corpus is mostly reporting the class prior.

\paragraph{Threshold selection is part of the model.} On the hate corpus the out-of-fold threshold is worth up to 0.087 F1 for classical models and 0.144 for the networks, more than the gap between the best and the median family. A comparison that does not tune thresholds per model is not comparing the models' best operating points.

\paragraph{The best cross-validated model is not always the best test model, and that is fine.} The CNN was selected on the hate corpus by cross-validated F1 and then scored 0.064 below its cross-validated mean on the test split, where calibrated classical models scored within 0.02 of theirs. The protocol reports this honestly instead of selecting on the test split, which would have produced a higher number and a biased one. The mechanism, an over-confident score scale that moves with early stopping, is identified and has a known remedy.

\paragraph{Representations interact with the target.} The ordering of representations is stable (sparse $n$-grams first, Doc2Vec last), but the size of the gaps depends on how the labels were made: on the lexicon-labelled vaccine corpus, where the target is a function of specific words, count features beat embeddings by 0.2 macro F1. Benchmarks that use weak labels should say so, because they measure the representations' ability to reconstruct the labeller.

\paragraph{Text length changes the winner.} On short documents the networks and tree ensembles are nominally ahead; on long balanced reviews the linear models on TF-IDF are the ceiling of the grid and the networks trained from scratch do not reach them. Only pretraining changes this picture, at a cost that is two orders of magnitude higher.

\paragraph{The 2022 study was directionally right and numerically wrong.} Its ranking of the classical models survives; its absolute numbers do not, and its headline number was an artefact. That is probably true of a large share of applied comparison studies, which is the strongest argument for protocols like the one defined here.

\section{Threats to validity}
\paragraph{Construct validity.} The vaccine corpus is labelled by a lexicon; its scores measure agreement with TextBlob, not sentiment as a human would judge it. The hate corpus has human labels from a competition without published guidelines, and mixes overt abuse with political content labelled by hashtag; its ceiling reflects label noise as much as model limits. IMDB has clean human labels derived from ratings.

\paragraph{Internal validity.} Test scores come from one stratified holdout per Twitter corpus and from the official IMDB split. Cross-validated means and fold standard deviations are reported alongside, and no claim in this thesis rests on a difference smaller than one fold standard deviation. Thresholds are chosen on fold models and applied to refitted ones; the cost of that transfer is measured, not assumed away. Hyper-parameters were set once from common practice rather than searched per cell; a per-cell search would raise most scores a little and would require a second level of nesting to remain unbiased.

\paragraph{External validity.} Both Twitter corpora date from 2016 to 2021; vocabulary, hashtags and platform norms have moved since, and a model trained on them would drift, which is what the released monitor measures. The grid excludes pretrained language models, so its conclusions concern models trained from scratch; the IMDB comparison with published results bounds the gap. Only English is considered.

\paragraph{Reproducibility.} All code, configuration and figures are public; corpora are obtained from their original sources. Seeds are fixed, but Gensim's multithreaded training makes the embedding-based cells reproducible only to about $10^{-3}$ in F1, and PyTorch on different hardware may differ at the same order. Timings are from one laptop.

\section{Lessons for applied practice}
Three habits would have prevented every defect of the 2022 study and cost nothing: put every fitted transformation inside the cross-validation loop; never select on the split you report; and tune the threshold on out-of-fold predictions while reporting precision-recall alongside ROC. A fourth habit, keeping the experiment in configuration and the outputs regenerable, turns a one-off comparison into an asset that can be re-run when a library, a corpus or a question changes. The system described in Chapter~\ref{ch:system} is one implementation of those habits; the habits matter more than the implementation.
